% ===========================================================================
% Styles TikZ partagés par toutes les vues d'architecture.
% Définis une seule fois ici : deux diagrammes du même rapport ne doivent pas
% employer deux vocabulaires graphiques différents.
%
% Convention de couleur — chaque teinte porte un sens, et un seul :
%   piqbitInk     : élément du système PIQBIT (ce que le projet a construit)
%   piqbitExt     : dépendance externe (toutes facultatives, avec repli)
%   piqbitActor   : acteur humain
%   piqbitStore   : donnée persistante
%   piqbitWarn    : frontière ou constat à signaler honnêtement
%
% ATTENTION : aucune ligne vide à l'intérieur de \tikzset{...}. Une ligne vide
% y vaut un \par, que pgfkeys refuse — l'erreur remontée est obscure
% (« Paragraph ended before \pgfkeys@addpath was complete »).
% ===========================================================================
% Même raison que pour la colonne L de config/commands.tex : le texte d'un
% nœud à `text width` est au fer à gauche, donc TeX le juge systématiquement
% « underfull » en le comparant à une ligne justifiée. L'avertissement est
% structurel et ne signale rien ; on le neutralise à l'échelle de l'image,
% pas du document.
\tikzset{every picture/.append style={execute at begin picture={\hbadness=10000\relax}}}

\definecolor{piqbitInk}{HTML}{06222A}   % teal profond — identité du produit
\definecolor{piqbitInkLite}{HTML}{E7EEF0}
\definecolor{piqbitExt}{HTML}{6B7C85}   % gris froid — hors frontière
\definecolor{piqbitExtLite}{HTML}{F0F2F3}
\definecolor{piqbitActor}{HTML}{3D4F4F}
\definecolor{piqbitStore}{HTML}{1F5A6B}
\definecolor{piqbitStoreLite}{HTML}{E4EEF1}
\definecolor{piqbitWarn}{HTML}{A8582F}  % ocre — jamais décoratif
\definecolor{piqbitWarnLite}{HTML}{FAF0E9}
%
\tikzset{
  %% ─── Blocs ─────────────────────────────────────────────────────────────
  pqbox/.style={
    draw=piqbitInk, fill=piqbitInkLite, line width=0.7pt,
    rounded corners=2pt, inner sep=5pt, align=center,
    font=\footnotesize, text=piqbitInk, minimum height=8mm
  },
  pqcore/.style={
    pqbox, fill=piqbitInk, text=white, line width=0.9pt,
    font=\footnotesize\bfseries
  },
  pqext/.style={
    draw=piqbitExt, fill=piqbitExtLite, line width=0.7pt, dashed,
    rounded corners=2pt, inner sep=5pt, align=center,
    font=\footnotesize, text=piqbitExt, minimum height=8mm
  },
  pqactor/.style={
    draw=piqbitActor, fill=white, line width=0.7pt,
    rounded corners=8pt, inner sep=5pt, align=center,
    font=\footnotesize, text=piqbitActor, minimum height=8mm
  },
  pqstore/.style={
    draw=piqbitStore, fill=piqbitStoreLite, line width=0.7pt,
    cylinder, shape border rotate=90, aspect=0.22,
    inner sep=4pt, align=center, font=\scriptsize, text=piqbitStore,
    minimum height=11mm, minimum width=19mm
  },
  pqwarn/.style={
    draw=piqbitWarn, fill=piqbitWarnLite, line width=0.8pt,
    rounded corners=2pt, inner sep=5pt, align=center,
    font=\scriptsize, text=piqbitWarn
  },
  pqnote/.style={
    font=\scriptsize\itshape, text=piqbitExt, align=left, inner sep=2pt
  },
  %% ─── Liens ─────────────────────────────────────────────────────────────
  pqflow/.style={
    -{Stealth[length=2.2mm]}, draw=piqbitInk, line width=0.7pt
  },
  pqflowback/.style={
    {Stealth[length=2.2mm]}-, draw=piqbitInk, line width=0.7pt
  },
  pqbi/.style={
    {Stealth[length=2.2mm]}-{Stealth[length=2.2mm]}, draw=piqbitInk,
    line width=0.7pt
  },
  %% Dépendance facultative : toujours en pointillés. Le trait plein est
  %% réservé à ce dont le système ne peut pas se passer.
  pqopt/.style={
    -{Stealth[length=2.2mm]}, draw=piqbitExt, line width=0.7pt, dashed
  },
  pqlbl/.style={
    font=\scriptsize, text=piqbitActor, fill=white, inner sep=1.5pt
  },
  %% ─── Frontières ────────────────────────────────────────────────────────
  pqboundary/.style={
    draw=piqbitInk, line width=0.9pt, dash pattern=on 3pt off 2pt,
    rounded corners=4pt, inner sep=7pt
  },
  pqtrust/.style={
    draw=piqbitWarn, line width=0.9pt, dash pattern=on 3pt off 2pt,
    rounded corners=4pt, inner sep=7pt
  },
  pqgroup/.style={
    draw=piqbitExt, line width=0.6pt, dotted,
    rounded corners=3pt, inner sep=6pt
  },
  %% `fill=white` n'est pas décoratif : le titre d'un groupe se pose SUR le
  %% cadre en pointillés qu'il nomme, et le trait le barrait.
  pqcaption/.style={
    font=\scriptsize\bfseries, text=piqbitInk, fill=white, inner sep=2pt
  },
  pqcaptionwarn/.style={
    font=\scriptsize\bfseries, text=piqbitWarn, fill=white, inner sep=2pt
  },
  %% ─── Diagrammes de séquence ────────────────────────────────────────────
  pqlifeline/.style={
    draw=piqbitExt, line width=0.5pt, dash pattern=on 2pt off 2pt
  },
  pqhead/.style={
    draw=piqbitInk, fill=piqbitInkLite, line width=0.7pt,
    inner xsep=4pt, inner ysep=3pt, align=center,
    font=\scriptsize\bfseries, text=piqbitInk
  },
  pqactivation/.style={
    draw=piqbitInk, fill=piqbitInk!12, line width=0.5pt
  },
  pqmsg/.style={
    -{Stealth[length=1.8mm]}, draw=piqbitInk, line width=0.6pt,
    font=\scriptsize
  },
  pqreturn/.style={
    -{Stealth[length=1.8mm]}, draw=piqbitExt, line width=0.6pt, dashed,
    font=\scriptsize
  },
  pqmsglbl/.style={
    font=\scriptsize, text=piqbitActor, above, midway, inner sep=1.5pt
  },
  %% ─── Modèle de données ─────────────────────────────────────────────────
  pqentity/.style={
    draw=piqbitStore, fill=white, line width=0.7pt,
    rectangle split, rectangle split parts=2,
    rectangle split part fill={piqbitStoreLite, white},
    rounded corners=1pt, inner sep=3pt, align=left,
    font=\scriptsize, text=piqbitInk
  },
  pqrel/.style={
    draw=piqbitStore, line width=0.6pt
  },
  pqcard/.style={
    font=\tiny, text=piqbitStore, inner sep=1.5pt
  }
}
%
% Légende compacte réutilisable sous les vues d'architecture.
\newcommand{\pqlegend}[1]{%
  \par\smallskip\noindent{\scriptsize\textit{Légende — #1}}\par
}
